\documentclass[border=4pt]{standalone}
\usepackage{amsmath,amssymb,mathtools,xcolor,varwidth}
\definecolor{timeblue}{HTML}{1E6FE0}
\definecolor{axisorange}{HTML}{E8770A}
\definecolor{ratiopurple}{HTML}{8E24AA}
\begin{document}
\begin{varwidth}{28em}
\large\bfseries
Kepler had noticed a haunting regularity in the heavens: the farther a planet      orbits from the Sun, the longer its year — and the relationship is exact. Squaring      \textcolor{timeblue}{the periodic time $T$} always tracks \textcolor{axisorange}{the cube of the greater axis $a$},      so that \textcolor{timeblue}{$T^2$} $\propto$ \textcolor{axisorange}{$a^3$}. Equivalently, \textcolor{timeblue}{the period}      grows as \textcolor{ratiopurple}{the $3/2$ power} — what Newton calls the \textcolor{ratiopurple}{sesquiplicate ratio} —      of the axis. What Kepler found only by patient observation, Newton now \textcolor{axisorange}{proves} from      the inverse-square law: "the periodic times in ellipses are in the sesquiplicate ratio of their greater axes."
\end{varwidth}
\end{document}
